% ROUND6-REFEREE-POSITIVE-CLOSURE
\section{Round-six positive closure: observed posterior games and canonical statistical tangents}
\label{sec:r6-c1}

The exact finite information state is the complete resolved history together
with the posterior law of the microscopic hard-sphere state under the already
controlled law.  Every posterior transition contains two distinct steps:
controlled prediction and disintegration by the newly observed block.  The
expected posterior-error exponent is Chernoff/testing information, whereas the
typical log posterior odds have the directed relative-entropy rate.  Local
asymptotic normality is centred at the exact finite-volume mean.

\subsection{Resolved observations and posterior kernels}

Let \(Z_t^\varepsilon\) be the complete hard-sphere microstate and let
\(Y_k^\varepsilon\) be the graph-completed resolved observation on the block
\([t_k,t_{k+1}]\).  The observation may contain finitely many density and
actual-contact cylinders and is a Borel map of the block path.  Put
\[
\mathcal F_{t_k}^{\rm res}=\sigma(Y_0^\varepsilon,\ldots,Y_{k-1}^\varepsilon),
\qquad
\Pi_k^\varepsilon=
\operatorname{Law}(Z_{t_k}^\varepsilon\mid
\mathcal F_{t_k}^{\rm res}).
\]
Regular conditional laws exist because both microstate and graph-completed
observation spaces are standard Borel.

For compact controls \(u\in U\), \(v\in V\), let
\(A_k^{u,v}(Z_{[t_k,t_{k+1}]})\) be a bounded particle/contact work.  Given a
posterior \(\pi\), define the exact conditional log normalizer
\[
\Lambda_{k,\varepsilon}^{u,v}(\pi)=
\frac1{\mu_\varepsilon}\log
\int e^{\mu_\varepsilon A_k^{u,v}(Z)}\,
\mathbf P_{\Delta_k}^\varepsilon(\pi,dZ),
\]
where \(\mathbf P_{\Delta_k}^\varepsilon(\pi,\cdot)\) is the deterministic
block-path law obtained by starting the Liouville flow from \(\pi\).
The controlled predictive block law is
\[
\overline{\mathbf P}_{k}^{u,v}(\pi,dZ)=
 e^{\mu_\varepsilon(A_k^{u,v}(Z)-
 \Lambda_{k,\varepsilon}^{u,v}(\pi))}
\mathbf P_{\Delta_k}^\varepsilon(\pi,dZ).
\]
Let \(O_k(Z)=Y_k\) be the new observation.  Disintegrate
\(\overline{\mathbf P}_{k}^{u,v}\) with respect to \(O_k\) and denote the
conditional terminal microstate law by
\[
\mathcal B_{k,\varepsilon}^{u,v}(\pi,y;\,dZ').
\]

\begin{lemma}[Controlled prediction--observation Bayes kernel]
\label{lem:r6-c1-bayes}
For every posterior \(\pi\), the predictive density above has total mass one.
The map
\[
(\pi,u,v)\longmapsto
\overline{\mathbf P}_{k}^{u,v}(\pi,\cdot)
\]
is Borel, and a jointly Borel version of
\(\mathcal B_{k,\varepsilon}^{u,v}\) exists.  Under the controlled law the
posterior update is exactly
\[
\Pi_{k+1}^{\varepsilon,u,v}
=
\mathcal B_{k,\varepsilon}^{u_k,v_k}
(\Pi_k^{\varepsilon,u,v},Y_k^\varepsilon;\cdot).
\]
Thus the update includes conditioning on the new observation; it is not merely
the tilted predictive pushforward.
\end{lemma}

\begin{proof}
The definition of \(\Lambda_{k,\varepsilon}^{u,v}\) normalizes the predictive
law.  Boundedness and Borel dependence of the work, together with continuity
of deterministic pushforward on bounded cylinders, give Borel dependence of
the predictive kernel.  Standard Borel disintegration provides a jointly
measurable posterior kernel after modifying one null set for a countable
generating algebra.  Applying Bayes' formula first to the exponential block
change of law and then to the sigma-field generated by \(O_k\) gives the
stated update.
\end{proof}

\subsection{Strategy-dependent laws and the exact DPP}

A relaxed control is a probability on \(U\) or \(V\).  An Elliott--Kalton
strategy maps the opponent's past relaxed controls and the resolved history to
the next relaxed control, nonanticipatively.  The information state is
\[
S_k^\varepsilon=(H_k^\varepsilon,\Pi_k^\varepsilon),
\]
where \(H_k^\varepsilon=(Y_0^\varepsilon,\ldots,Y_{k-1}^\varepsilon)\).

\begin{theorem}[Canonical strategy-law construction]
\label{thm:r6-c1-law}
For every pair of nonanticipative relaxed strategies and every initial
posterior, the Bayes kernels of Lemma~\ref{lem:r6-c1-bayes} determine a unique
controlled probability law on histories, posteriors, controls, and
microstates.  Under that law each block likelihood increment is conditionally
mean one with respect to the already controlled posterior.  The family is
stable under concatenation at stopping grid times.
\end{theorem}

\begin{proof}
At the first block, evaluate the strategies at the initial history and use the
normalized predictive/observation kernel.  Inductively, the realized history
and posterior determine the next relaxed controls and hence the next kernel.
Ionescu--Tulcea gives a unique path law.  The normalizer is evaluated under the
posterior generated by all preceding controlled kernels and observations, so
the next likelihood increment has conditional mean one.  Disintegration at a
stopping grid time and restarting the same recursion proves concatenation.
\end{proof}

Let the terminal payoff be bounded and let the running reward be continuous in
the information state and relaxed controls.

\begin{theorem}[Exact posterior-state dynamic programming]
\label{thm:r6-c1-dpp}
The upper and lower finite-grid game values are Borel functions of
\((H_k,\Pi_k)\) and satisfy the exact dynamic programming recursion on that
state.  Compactness of relaxed control spaces gives measurable selectors.  A
time-zero preparation action, a reward-only feedback action, and an adaptive
law action define three different games:
\begin{enumerate}[label=(\roman*)]
\item a preparation action is chosen once and remains a static state
coordinate through all blocks;
\item reward-only feedback changes the running payoff but not the transition
kernel;
\item adaptive law control changes the normalized predictive kernel and hence
the posterior transition.
\end{enumerate}
No pointwise Isaacs equation is assigned to the first game unless the static
preparation action is retained in the state.
\end{theorem}

\begin{proof}
Condition the value at the next grid time.  The state
\((H_{k+1},\Pi_{k+1})\) is a measurable function of the current state, relaxed
controls, and new observation through
Lemma~\ref{lem:r6-c1-bayes}.  The tower property under the unique law of
Theorem~\ref{thm:r6-c1-law} gives the recursion.  Probability measures on
compact action sets are compact, and continuity of bounded work/reward makes
the one-step value lower/upper semicontinuous, so measurable minimax selectors
exist.  The three timing statements follow directly from which coordinate of
the kernel or payoff the action changes.
\end{proof}

\subsection{Posterior trace classes and the kinetic contraction}

Let
\[
\|\pi\|_{\rm tr}=\|G(\pi)\|_{\alpha,\beta}^{\rm tr}
\]
be the B2 trace-hierarchy norm.  Conditional expectation gives
\[
\mathbb E\|\Pi_k^\varepsilon\|_{\rm tr}
\le C_k\|\Pi_0^\varepsilon\|_{\rm tr}e^{Ck\Delta e^{CR}}
\]
for bounded control sources.  Thus posterior trace norms need not be uniformly
bounded for every rare observation, but they have exponential/entropy
containment under every strategy.

\begin{lemma}[Strategy-uniform conditional source charts]
\label{lem:r6-c1-source}
On each posterior trace ball, the next-block conditional particle/contact
pressure is the B2 pressure with that posterior hierarchy as initial state and
has the same analytic radius and derivative bounds, with constants depending
only on the ball.  Under bounded strategies, posterior paths leave an
increasing sequence of such balls with exponentially vanishing probability.
\end{lemma}

\begin{proof}
The B2 trace-class theorem is linear in the initial complete hierarchy and its
operator bounds depend only on its trace norm.  Conditioning changes the
initial hierarchy but not the local hard-sphere dynamics.  Hence the same
connected expansion applies on a trace ball.  The normalized block likelihood
and Bayes disintegration preserve the average trace norm; the entropy
inequality applied to its exponential Lyapunov truncations gives exponential
containment uniformly over bounded control works.  Iteration over finitely
many blocks gives the claim.
\end{proof}

Let \(q^{u,v}\) be the limiting positive bounded actual-contact multiplier
selected by the control.  Its kinetic Hamiltonian is
\[
\mathbb H^{u,v}(f,p)=
\langle f,v\cdot\nabla p\rangle+
\int q^{u,v}(e^{\Delta p}-1)dA_f+c(f,u,v).
\]

\begin{theorem}[Posterior-to-kinetic contraction]
\label{thm:r6-c1-contraction}
On every finite control grid, conditionally on the resolved density/contact
history, the posterior finite marginals satisfy a uniform conditional Laplace
principle whose unique regular minimizer is the controlled B2 biased hierarchy.
Consequently the posterior coordinate contracts, in finite correlation
projections, to that hierarchy and its one-particle density.  As the grid mesh
tends to zero, the game values converge to the unique B4 viscosity solution of
\[
\partial_tU+\sup_u\inf_v
\mathbb H^{u,v}
\left(f,\frac{\delta U}{\delta f}\right)=0
\]
when the Isaacs equality holds.  The exact finite information state remains
\((H,\Pi)\); density reduction is a proved kinetic limit, not a finite-scale
identity.
\end{theorem}

\begin{proof}
On one block, apply the B2 normalized bounded-source tilt starting from the
current posterior trace hierarchy.  The conditional pressure derivatives and
strict quotient covariance identify one regular biased hierarchy, and the
conditional large-deviation upper and lower bounds give exponential
concentration there.  Lemma~\ref{lem:r6-c1-source} makes the estimates uniform
on exponentially containing trace balls.  Induct over the finite grid,
carrying the posterior update with the new observations.  Finite correlation
projections then converge to the controlled hierarchy.

The discrete DPP is consistent with the displayed Hamiltonian by the B2
one-block pressure.  B4 compact containment and comparison identify the upper
and lower half-relaxed limits as the mesh vanishes.  Isaacs equality identifies
the two game values.
\end{proof}

\subsection{Typical odds and expected posterior error}

Let \(\{\mathbb P_{\theta,\varepsilon}:\theta\in\Theta\subset\mathbb R^d\}\)
be a compact canonical exponential path family with
\[
\frac{d\mathbb P_{\theta,\varepsilon}}
{d\mathbb P_{0,\varepsilon}}
=\exp\{\mu_\varepsilon(\theta\cdot X_\varepsilon-Q_\varepsilon(\theta))\}.
\]
Assume the prior has a continuous density bounded above and below near the
true interior parameter \(\theta_0\).  Posterior densities are always taken
with respect to Lebesgue measure or this prior; no ratio to the mass of a
singleton is used.

For \(\theta\ne\theta_0\), define the Chernoff information
\[
\mathcal C(\theta_0,\theta)=
-\inf_{0\le s\le1}
\left[
Q((1-s)\theta_0+s\theta)
-(1-s)Q(\theta_0)-sQ(\theta)
\right].
\]

\begin{theorem}[KL odds and Chernoff posterior-error exponents]
\label{thm:r6-c1-testing}
Under \(\mathbb P_{\theta_0,\varepsilon}\), the typical posterior log-density
ratio obeys
\[
\frac1{\mu_\varepsilon}\log
\frac{r_T^\varepsilon(\theta)}{r_T^\varepsilon(\theta_0)}
\longrightarrow
-\mathcal K(\theta_0\Vert\theta)
\]
in probability, locally uniformly away from nonidentifiable parameters.  For
every neighborhood \(O\) of \(\theta_0\),
\[
\limsup_{\varepsilon\to0}
\frac1{\mu_\varepsilon}\log
\mathbb E_{\theta_0}
\int_{O^c}r_T^\varepsilon(\theta)d\theta
\le-\inf_{\theta\in O^c}\mathcal C(\theta_0,\theta).
\]
The right side is strictly negative when the family is identifiable.
\end{theorem}

\begin{proof}
The log likelihood ratio divided by \(\mu_\varepsilon\) converges by the
pressure law of large numbers to its expectation, which is minus the directed
relative-entropy rate.  Prior-density ratios are subexponential, giving the
first assertion.

For expected posterior error, use
\(x/(1+x)\le x^s\), \(0\le s\le1\), for each likelihood ratio and optimize in
\(s\).  The exact exponential-family moment is the Chernoff pressure displayed
above.  Cover the compact complement of \(O\) by finitely many neighborhoods
on which pressure and prior density vary by at most \(o(1)\); analytic
compactness makes the cover independent of \(\varepsilon\).  Integrating the
pointwise bounds and optimizing proves the exponent.  Strict convexity on the
identifiable quotient makes the infimum positive.
\end{proof}

\subsection{Exact-centred LAN and Bernstein--von Mises}

Define
\[
\Delta_\varepsilon=\sqrt{\mu_\varepsilon}
\bigl(X_\varepsilon-DQ_\varepsilon(\theta_0)\bigr),
\qquad
I_\varepsilon=D^2Q_\varepsilon(\theta_0).
\]

\begin{theorem}[Finite-volume-centred canonical LAN]
\label{thm:r6-c1-lan}
Uniformly for \(h\) in compact subsets of \(\mathbb R^d\),
\[
\log\frac{d\mathbb P_{\theta_0+h/\sqrt{\mu_\varepsilon},\varepsilon}}
{d\mathbb P_{\theta_0,\varepsilon}}
=h\cdot\Delta_\varepsilon-\frac12h^TI_\varepsilon h+o_{\mathbb P}(1).
\]
The score is exactly centred for every finite \(\varepsilon\).  Moreover
\(\Delta_\varepsilon\Rightarrow N(0,I)\), where
\(I=D^2Q(\theta_0)>0\) on the identifiable quotient.
\end{theorem}

\begin{proof}
Taylor-expand the exact finite pressure at
\(\theta_0+h/\sqrt{\mu_\varepsilon}\).  The linear term is
\(DQ_\varepsilon(\theta_0)h/\sqrt{\mu_\varepsilon}\) and cancels exactly with
the chosen score centre.  Uniform third Cauchy bounds on a larger complex
chart make the remainder \(O(|h|^3/\sqrt{\mu_\varepsilon})\).  The B3
exact-centred Gaussian theorem gives the score limit and pressure derivative
convergence gives \(I_\varepsilon\to I\).
\end{proof}

\begin{theorem}[Canonical Bernstein--von Mises theorem]
\label{thm:r6-c1-bvm}
Assume the canonical chart is finite dimensional, identifiable, the prior
density is positive and continuous at \(\theta_0\), and the compact
complements admit the Chernoff tests of
Theorem~\ref{thm:r6-c1-testing}.  Then the posterior law of
\(h=\sqrt{\mu_\varepsilon}(\theta-\theta_0)\), after centring by
\(I^{-1}\Delta_\varepsilon\), converges in total variation to
\(N(0,I^{-1})\).
\end{theorem}

\begin{proof}
The Chernoff bound makes posterior mass outside every fixed neighborhood
exponentially small.  On balls \(|h|\le M_\varepsilon\) with
\(M_\varepsilon\uparrow\infty\) and
\(M_\varepsilon^3/\sqrt{\mu_\varepsilon}\to0\), the LAN remainder is uniform.
The prior ratio converges uniformly to one.  Completing the square in the LAN
likelihood gives the Gaussian density with centre
\(I^{-1}\Delta_\varepsilon\); dominated convergence gives total variation on
the growing ball.  Positive definiteness gives Gaussian tail control, and the
Chernoff tests control the complement.
\end{proof}

\subsection{Smooth and discrete saddle envelopes}

On a smooth finite-dimensional preparation manifold, strict signed curvature
of the preparation cost plus the B3 quotient covariance gives an invertible
saddle Hessian.  The envelope Hessian is the corresponding Schur complement.
For finite action sets the value is an active-branch directional envelope and
no smooth Hessian is asserted at branch crossings.

\begin{theorem}[Round-six C1 closure]
\label{thm:r6-c1-main}
The finite hard-sphere control problem has an exact posterior-history
information state, normalized strategy-dependent laws, and an observed Bayes
DPP.  Its adaptive kinetic limit is the Isaacs semigroup obtained from the B2
controlled actual-contact pressure and B4 comparison.  Typical posterior odds
have the KL rate, expected posterior error has the Chernoff rate, and the
canonical statistical experiment satisfies exact-centred LAN and the stated
Bernstein--von Mises theorem.
\end{theorem}

\begin{proof}
The exact game follows from Theorems~\ref{thm:r6-c1-law} and
\ref{thm:r6-c1-dpp}; its kinetic limit is
Theorem~\ref{thm:r6-c1-contraction}.  The statistical statements are
Theorems~\ref{thm:r6-c1-testing}, \ref{thm:r6-c1-lan}, and
\ref{thm:r6-c1-bvm}.  The timing distinctions and saddle calculus were proved
on their respective exact state spaces above.
\end{proof}
